% solutions/obs-marginals-not-joint.tex
% Standalone proof dossier for the fixed-marginal A3 obstruction.
\documentclass[../main.tex]{subfiles}
\begin{document}

% === SOLUTION HEADER (proof provenance) =====================================
%   ledger-node : obs:marginals-not-joint
%   refines     : prop:a3-marginals-not-joint
%   bounded_by  : none
%   evidence_run: none
%   checked_by  : agent
%   author      : /root/a3
%   reviewer    : /root/audit_latex
%   review      : research/reviews/2026-08-22-a3-marginals-independent-audit.md
%   review level: independent agent; no human or Lean review is claimed
%   date        : 2026-08-22
% ============================================================================

\section*{Solution: fixed marginals do not control the joint weighted gap}
\label{sec:sol-obs-marginals-not-joint}

\begin{proposition}
\label{prop:sol-a3-marginals-not-joint}
For every $0<\varepsilon<1$ there is a positive, piecewise-smooth density $q_\varepsilon$ on
$\mathbb R^3$ whose three coordinate marginals are standard Cauchy and whose optimal constant in
\[
 \Var_{q_\varepsilon}(f)
 \le C_\varepsilon
 \int\sum_{j=1}^3(1+x_j^2)|\partial_jf|^2q_\varepsilon(x)\,\dd x
\]
satisfies $C_\varepsilon\asymp\varepsilon^{-1}$. In particular, fixed one-dimensional marginals
and their Hardy constants do not uniformly control the dependent joint constant. The same holds
in every dimension $d\ge3$ after adjoining independent Cauchy coordinates.
\end{proposition}

\begin{proof}
Let $p(x)=\{\pi(1+x^2)\}^{-1}$ have CDF $F$, put $A=(0,1/2)$ and $D=(1/2,1)$, and define on
$(0,1)^3$
\[
 c_\varepsilon(u)
 =4(1-\varepsilon)\bigl(\mathbf1_{A^3}(u)+\mathbf1_{D^3}(u)\bigr)+\varepsilon.
\]
It is positive and normalized because the two distinguished cubes have total volume $1/4$.
For any fixed value of one coordinate, integration over the other two gives
$4(1-\varepsilon)(1/2)^2+\varepsilon=1$. Thus every marginal of $c_\varepsilon$ is uniform, and
\[
 q_\varepsilon(x)
 =\prod_{j=1}^3p(x_j)\,
   c_\varepsilon(F(x_1),F(x_2),F(x_3))
\]
has three standard-Cauchy marginals.

Write $o=(1/2,1/2,1/2)$. Choose a smooth odd function $g$ on $S^2$ that is $1$ on the
all-negative spherical octant, $-1$ on the all-positive octant, and varies only through the six
mixed-sign octants. Set
\[
 \varphi(o+r\omega)=g(\omega)
\]
away from $o$, with an arbitrary value at $o$. Since $|\nabla\varphi|=O(r^{-1})$,
\[
 \int_0^\delta r^{-2}r^2\,\dd r<\infty,
\]
so $\varphi\in H^1((0,1)^3)$. Reflection $u\mapsto1-u$ preserves $c_\varepsilon$ and negates
$\varphi$, hence its mean is zero. The two plateau cubes have total
$c_\varepsilon$-mass $1-3\varepsilon/4$, and therefore
\[
 \Var_{c_\varepsilon}(\varphi)\ge1-\frac{3\varepsilon}{4}.
\]

For $f(x)=\varphi(F(x_1),F(x_2),F(x_3))$, the change of variables $u_j=F(x_j)$ gives
\begin{align*}
 &\int\sum_{j=1}^3(1+x_j^2)|\partial_jf|^2q_\varepsilon(x)\,\dd x\\
 &\quad=\int c_\varepsilon(u)\sum_{j=1}^3
 (1+F^{-1}(u_j)^2)p(F^{-1}(u_j))^2|\partial_j\varphi|^2\,\dd u.
\end{align*}
Now $(1+x^2)p(x)^2\le\pi^{-2}$, and $\nabla\varphi$ is supported in the mixed octants, where
$c_\varepsilon=\varepsilon$. Hence the energy is at most
\[
 \frac{\varepsilon}{\pi^2}\int|\nabla\varphi|^2\,\dd u,
\]
which yields $C_\varepsilon\ge c/\varepsilon$. Smooth $H^1$ approximants give the same lower
bound: choose their squared $H^1$ error to be $o(\varepsilon)$; the transformed energy
coefficient is bounded by $4/\pi^2$.

For completeness, $\varepsilon\le c_\varepsilon\le4-3\varepsilon$. Density comparison with the
product Cauchy law, whose product weighted constant is $4$, gives
$C_\varepsilon\le4(4-3\varepsilon)/\varepsilon$. Thus
$C_\varepsilon\asymp\varepsilon^{-1}$.
\end{proof}

\begin{corollary}[Coefficient-only energy cannot control nondegenerate scales]
Let $\pi$ be a joint coefficient--scale law whose scale marginal is nondegenerate. A Dirichlet
form differentiating only in the coefficient coordinates cannot satisfy a finite Poincar\'e
inequality on the joint space.
\end{corollary}

\begin{proof}
Choose a bounded, nonconstant function of the scale coordinates alone. It has positive variance
and zero coefficient-gradient energy.
\end{proof}

\end{document}
